\subsection{Properties of induced measures}

PROPOSITION 3. --- Let X be a locally compact subspace of T, and \(\lambda\) a complex measure on X. The following properties are equivalent:

\begin{enumerate}
\def\labelenumi{\alph{enumi})}
\item
  the canonical injection \(i:\mathrm{X}\to \mathrm{T}\) is \(\lambda\) -proper;
\item
  for every compact subset K of T, \(\mathbf{K} \cap \mathbf{X}\) is essentially \(\lambda\) -integrable;
\item
  every point \(t \in \mathbf{T}\) admits a neighborhood \(\mathbf{V}\) such that \(\mathbf{V} \cap \mathbf{X}\) is essentially \(\lambda\) -integrable;
\item
  there exists a measure \(\theta\) on T such that \(\theta_{\mathrm{X}} = \lambda\) .
\end{enumerate}

If these equivalent conditions are satisfied, we have, with notations as in d),

\[
\left(i (\lambda)\right) _ {\mathrm{X}} = \lambda \quad \text{and} \quad i (\lambda) = i (\theta_ {\mathrm{X}}) = \varphi_ {\mathrm{X}} \cdot \theta .\tag{7}
\]

The injection \(i\) being continuous, the equivalence of the properties a), b) and c) follows from Prop. 1 of §6, and the remark that follows it, applied to the positive measure \(|\lambda|\) . If \(\lambda\) is induced on X by a measure \(\theta\) on T, then \(|\lambda| = |\theta|_{\mathrm{X}}\) (formula (6)), consequently \(|\lambda| (\mathrm{K} \cap \mathrm{X}) = |\theta| (\mathrm{K} \cap \mathrm{X}) \leqslant |\theta| (\mathrm{K}) < +\infty\) (Prop. 1) for every compact subset K of T, so that d) implies b). Suppose that a) is satisfied, and let us show that \(\left(i(\lambda)\right)_{\mathrm{X}} = \lambda\) , which will imply d). Let \(g\) be an element of \(\mathcal{H}(\mathrm{X}; \mathbf{C})\) ; denoting by \(g'\) the extension by 0 of \(g\) to T we have, by the definition of induced measure and then by Prop. 7 of §6, No.~4.

\[
\int g d (i (\lambda)) _ {X} = \int g ^ {\prime} d (i (\lambda)) = \int (g ^ {\prime} \circ i) d \lambda = \int g d \lambda .
\]

This completes the proof of the equivalence of the four properties. If \(\lambda = \theta_{\mathrm{X}}\) and \(g\in \mathcal{K}(\mathrm{T};\mathbf{C})\) , then

\[
\int g d \bigl (i (\theta_ {\mathrm{X}}) \bigr) = \int (g \circ i) d (\theta_ {\mathrm{X}}) = \int g \varphi_ {\mathrm{X}} d \theta ,
\]

because \(g\varphi_{\mathbf{X}}\) is the extension by 0 of \(g \circ i\) to \(\mathbf{T}\) . This proves the second formula of (7).

COROLLARY 1. --- If X is closed, then every complex measure \(\lambda\) on X is induced by a measure on T.

For, if K is a compact set in T then \(K \cap X\) is compact, hence \(\lambda\) -integrable.

COROLLARY 2. --- Let \(\theta\) be a complex measure on T, \(\pi\) a \(\theta\) -proper mapping of T into a locally compact space Y, and \(\pi_{X}\) its restriction to X. Then \(\pi_{X}\) is \(\theta_{X}\) -proper, and \(\pi_{X}(\theta_{X}) = \pi(\varphi_{X} \cdot \theta)\) .

For, \(\pi_{X} = \pi \circ i\) , where i is the canonical injection \(X \to T\) . When \(\theta\) is positive the corollary may therefore be deduced from Prop. 3 and the transitivity of image measures ( \(\S6\) , No.~3, Prop. 4). The case of a non-positive complex measure then follows by linearity.

PROPOSITION 4. --- Let X and Y be two locally compact subspaces of T such that \(Y \subset X\) . If \(\theta\) is a complex measure on T, then the measure \((\theta_{\mathrm{X}})_{\mathrm{Y}}\) induced by \(\theta_{X}\) on Y is equal to \(\theta_{Y}\) (`transitivity of induced measures').

It suffices to observe that if g is an element of \(\mathcal{K}(\mathrm{Y};\mathbf{C})\) , then the extension by 0 of g to T may be obtained by extending by 0 the extension by 0 of g to X, or again, making use of the identifications of the Scholium, that \(\varphi_{\mathrm{Y}}\cdot\theta=\varphi_{\mathrm{Y}}(\varphi_{\mathrm{X}}\cdot\theta)\) (§5, No.~4, Prop. 8).

PROPOSITION 5. --- Let \((\lambda_{\alpha})_{\alpha\in A}\) be an increasing directed family of positive measures on T, admitting a supremum \(\lambda\) , and let X be a locally compact subspace of T. The family of induced measures \(\lambda_{\alpha}|X\) is then bounded above in \(\mathcal{M}(X)\) , and

\[
\sup _ {\alpha \in \mathrm{A}} (\lambda_ {\alpha} | \mathrm{X}) = \lambda | \mathrm{X}.\tag{8}
\]

In view of the identifications in the Scholium, this proposition is a special case of Prop. 5 of §5, No.~4.

COROLLARY. --- Let \((\mu_{i})_{i\in\mathbb{I}}\) be a summable family of positive measures on T, with sum \(\mu\) . The family of induced measures \(\mu_{i}|X\) is then summable, and

\[
\sum_ {i \in I} (\mu_ {i} | X) = \mu | X.\tag{9}
\]

PROPOSITION 6. --- Let \(\Lambda: t \mapsto \lambda_t\) be a \(\mu\) -adequate mapping of T into \(\mathcal{M}_+(\mathrm{X})\) , where X is a locally compact space that is countable at infinity, and let Y be a locally compact subspace of X. Set \(\int \lambda_t d\mu(t) = \nu\) . The mapping \(t \mapsto \lambda_t|\mathrm{Y}\) of T into \(\mathcal{M}_+(\mathrm{Y})\) is then \(\mu\) -adequate, and

\[
\int (\lambda_ {t} | \mathrm{Y}) d \mu (t) = \nu | \mathrm{Y}.\tag{10}
\]

Taking into account the identifications in the Scholium, this proposition is a special case of Prop. 7 of §5, No.~4.

